%
\begin{isabellebody}%
\setisabellecontext{Week{\isacharunderscore}{\kern0pt}{\isadigit{1}}{\isacharunderscore}{\kern0pt}sol}%
%
\isadelimtheory
\isanewline
%
\endisadelimtheory
%
\isatagtheory
\isacommand{theory}\isamarkupfalse%
\ Week{\isacharunderscore}{\kern0pt}{\isadigit{1}}{\isacharunderscore}{\kern0pt}sol\isanewline
\isakeyword{imports}\ Main\isanewline
\isakeyword{begin}%
\endisatagtheory
{\isafoldtheory}%
%
\isadelimtheory
%
\endisadelimtheory
%
\begin{isamarkuptext}%
\vspace{15ex}\ExerciseSheet{6}{}%
\end{isamarkuptext}\isamarkuptrue%
%
\begin{isamarkuptext}%
Important Note: Please download Isabelle from the link provided in the slides and bring a 
      laptop to the large group tutorial. Performing the proofs in Isabelle will be an integral part
      of those tutorials.%
\end{isamarkuptext}\isamarkuptrue%
%
\begin{isamarkuptext}%
Before you start, please open a new Isabelle/HOL theory file called Week\_1.thy and add the
       following to its beginning.

\noindent theory Week\_1\\
imports Main\\
begin%
\end{isamarkuptext}\isamarkuptrue%
%
\begin{isamarkuptext}%
\Exercise{Addition Commutes}%
\end{isamarkuptext}\isamarkuptrue%
%
\begin{isamarkuptext}%
Consider the function \isa{{\isacharparenleft}{\kern0pt}{\isacharplus}{\kern0pt}{\isacharparenright}{\kern0pt}} defined in Isabelle/HOL for natural numbers.

Prove the following theorem, first pen-and-paper, and then formally in Isabelle.%
\end{isamarkuptext}\isamarkuptrue%
\isacommand{lemma}\isamarkupfalse%
\ add{\isacharunderscore}{\kern0pt}commutes{\isacharcolon}{\kern0pt}\ {\isachardoublequoteopen}{\isacharparenleft}{\kern0pt}n{\isacharcolon}{\kern0pt}{\isacharcolon}{\kern0pt}nat{\isacharparenright}{\kern0pt}\ {\isacharplus}{\kern0pt}\ m\ {\isacharequal}{\kern0pt}\ m\ {\isacharplus}{\kern0pt}\ n{\isachardoublequoteclose}%
\begin{isamarkuptext}%
\paragraph{Hints:} 

 1. Perform the proof by inuction on \isa{n}, using substitutions \isa{subst}, and forward reasoning using
   the method \isa{rule}.

 2. You will need to identify multiple lemmas about how \isa{add} works. Do not prove those lemmas. Use 
    \isa{sorry} as the proof. This is the top-down approach. 

 3. There is a proof which uses the following lemmas:

   \isa{add{\isacharunderscore}{\kern0pt}{\isadigit{0}}{\isacharunderscore}{\kern0pt}right{\isacharcomma}{\kern0pt}\ add{\isacharunderscore}{\kern0pt}{\isadigit{0}}{\isacharunderscore}{\kern0pt}left{\isacharcomma}{\kern0pt}\ add{\isacharunderscore}{\kern0pt}Suc{\isacharunderscore}{\kern0pt}right{\isacharcomma}{\kern0pt}\ add{\isacharunderscore}{\kern0pt}Suc}%
\end{isamarkuptext}\isamarkuptrue%
%
\isadelimproof
%
\endisadelimproof
%
\isatagproof
\isacommand{proof}\isamarkupfalse%
{\isacharparenleft}{\kern0pt}induction\ m{\isacharparenright}{\kern0pt}\isanewline
\ \ \isacommand{case}\isamarkupfalse%
\ {\isadigit{0}}\isanewline
\ \ \isacommand{have}\isamarkupfalse%
\ {\isadigit{2}}{\isacharcolon}{\kern0pt}\ {\isachardoublequoteopen}n\ {\isacharplus}{\kern0pt}\ {\isadigit{0}}\ {\isacharequal}{\kern0pt}\ n{\isachardoublequoteclose}\isanewline
\ \ \ \ \isacommand{apply}\isamarkupfalse%
{\isacharparenleft}{\kern0pt}subst\ {\isacharparenleft}{\kern0pt}{\isadigit{1}}{\isacharparenright}{\kern0pt}\ add{\isacharunderscore}{\kern0pt}{\isadigit{0}}{\isacharunderscore}{\kern0pt}right{\isacharparenright}{\kern0pt}\isanewline
\ \ \ \ \isacommand{{\isachardot}{\kern0pt}{\isachardot}{\kern0pt}}\isamarkupfalse%
\isanewline
\ \ \isacommand{have}\isamarkupfalse%
\ {\isadigit{4}}{\isacharcolon}{\kern0pt}\ {\isachardoublequoteopen}{\isadigit{0}}\ {\isacharplus}{\kern0pt}\ n\ {\isacharequal}{\kern0pt}\ n{\isachardoublequoteclose}\isanewline
\ \ \ \ \isacommand{apply}\isamarkupfalse%
{\isacharparenleft}{\kern0pt}subst\ {\isacharparenleft}{\kern0pt}{\isadigit{1}}{\isacharparenright}{\kern0pt}\ add{\isacharunderscore}{\kern0pt}{\isadigit{0}}{\isacharunderscore}{\kern0pt}left{\isacharparenright}{\kern0pt}\isanewline
\ \ \ \ \isacommand{{\isachardot}{\kern0pt}{\isachardot}{\kern0pt}}\isamarkupfalse%
\isanewline
\isanewline
\ \ \isacommand{have}\isamarkupfalse%
\ {\isadigit{5}}{\isacharcolon}{\kern0pt}\ {\isachardoublequoteopen}n\ {\isacharplus}{\kern0pt}\ {\isadigit{0}}\ {\isacharequal}{\kern0pt}\ n{\isachardoublequoteclose}\isanewline
\ \ \ \ \isacommand{using}\isamarkupfalse%
\ {\isadigit{2}}\isanewline
\ \ \ \ \isacommand{{\isachardot}{\kern0pt}}\isamarkupfalse%
\isanewline
\ \ \isacommand{from}\isamarkupfalse%
\ {\isadigit{5}}\ \isacommand{show}\isamarkupfalse%
\ {\isacharquery}{\kern0pt}case\isanewline
\ \ \ \ \isacommand{by}\isamarkupfalse%
{\isacharparenleft}{\kern0pt}subst\ {\isadigit{4}}{\isacharparenright}{\kern0pt}\isanewline
\isacommand{next}\isamarkupfalse%
\isanewline
\ \ \isacommand{case}\isamarkupfalse%
\ {\isacharparenleft}{\kern0pt}Suc\ m{\isacharparenright}{\kern0pt}\isanewline
\isanewline
\ \ \isacommand{hence}\isamarkupfalse%
\ {\isadigit{2}}{\isacharcolon}{\kern0pt}\ {\isachardoublequoteopen}n\ {\isacharplus}{\kern0pt}\ {\isacharparenleft}{\kern0pt}Suc\ m{\isacharparenright}{\kern0pt}\ {\isacharequal}{\kern0pt}\ Suc\ {\isacharparenleft}{\kern0pt}n\ {\isacharplus}{\kern0pt}\ m{\isacharparenright}{\kern0pt}{\isachardoublequoteclose}\isanewline
\ \ \ \ \isacommand{apply}\isamarkupfalse%
{\isacharparenleft}{\kern0pt}subst\ add{\isacharunderscore}{\kern0pt}Suc{\isacharunderscore}{\kern0pt}right{\isacharparenright}{\kern0pt}\isanewline
\ \ \ \ \isacommand{{\isachardot}{\kern0pt}{\isachardot}{\kern0pt}}\isamarkupfalse%
\isanewline
\isanewline
\ \ \isacommand{hence}\isamarkupfalse%
\ {\isadigit{4}}{\isacharcolon}{\kern0pt}\ {\isachardoublequoteopen}{\isacharparenleft}{\kern0pt}Suc\ m{\isacharparenright}{\kern0pt}\ {\isacharplus}{\kern0pt}\ n\ {\isacharequal}{\kern0pt}\ Suc\ {\isacharparenleft}{\kern0pt}m\ {\isacharplus}{\kern0pt}\ n{\isacharparenright}{\kern0pt}{\isachardoublequoteclose}\isanewline
\ \ \ \ \isacommand{apply}\isamarkupfalse%
{\isacharparenleft}{\kern0pt}subst\ add{\isacharunderscore}{\kern0pt}Suc{\isacharparenright}{\kern0pt}\isanewline
\ \ \ \ \isacommand{{\isachardot}{\kern0pt}{\isachardot}{\kern0pt}}\isamarkupfalse%
\isanewline
\isanewline
\ \ \isacommand{hence}\isamarkupfalse%
\ {\isadigit{6}}{\isacharcolon}{\kern0pt}\ {\isachardoublequoteopen}Suc\ {\isacharparenleft}{\kern0pt}n\ {\isacharplus}{\kern0pt}\ m{\isacharparenright}{\kern0pt}\ {\isacharequal}{\kern0pt}\ Suc\ {\isacharparenleft}{\kern0pt}m\ {\isacharplus}{\kern0pt}\ n{\isacharparenright}{\kern0pt}{\isachardoublequoteclose}\isanewline
\ \ \ \ \isanewline
\ \ \ \ \isacommand{thm}\isamarkupfalse%
\ Suc{\isachardot}{\kern0pt}IH\isanewline
\ \ \ \ \isanewline
\ \ \ \ \isacommand{apply}\isamarkupfalse%
{\isacharparenleft}{\kern0pt}subst\ Suc{\isachardot}{\kern0pt}IH{\isacharparenright}{\kern0pt}\isanewline
\ \ \ \ \isacommand{{\isachardot}{\kern0pt}{\isachardot}{\kern0pt}}\isamarkupfalse%
\isanewline
\isanewline
\ \ \isacommand{hence}\isamarkupfalse%
\ {\isadigit{7}}{\isacharcolon}{\kern0pt}\ {\isachardoublequoteopen}n\ {\isacharplus}{\kern0pt}\ {\isacharparenleft}{\kern0pt}Suc\ m{\isacharparenright}{\kern0pt}\ {\isacharequal}{\kern0pt}\ Suc\ {\isacharparenleft}{\kern0pt}m\ {\isacharplus}{\kern0pt}\ n{\isacharparenright}{\kern0pt}{\isachardoublequoteclose}\isanewline
\ \ \ \ \isacommand{by}\isamarkupfalse%
{\isacharparenleft}{\kern0pt}subst\ {\isacharparenleft}{\kern0pt}asm{\isacharparenright}{\kern0pt}\ {\isadigit{2}}{\isacharbrackleft}{\kern0pt}symmetric{\isacharbrackright}{\kern0pt}{\isacharparenright}{\kern0pt}\isanewline
\isanewline
\ \ \isacommand{thus}\isamarkupfalse%
\ {\isacharquery}{\kern0pt}case\isanewline
\ \ \ \ \isacommand{by}\isamarkupfalse%
{\isacharparenleft}{\kern0pt}subst\ {\isacharparenleft}{\kern0pt}asm{\isacharparenright}{\kern0pt}\ {\isadigit{4}}{\isacharbrackleft}{\kern0pt}symmetric{\isacharbrackright}{\kern0pt}{\isacharparenright}{\kern0pt}\isanewline
\isanewline
\isacommand{qed}\isamarkupfalse%
%
\endisatagproof
{\isafoldproof}%
%
\isadelimproof
%
\endisadelimproof
%
\begin{isamarkuptext}%
\Exercise{Multiplication is Monotone}%
\end{isamarkuptext}\isamarkuptrue%
%
\begin{isamarkuptext}%
Consider the multiplication function \isa{{\isacharparenleft}{\kern0pt}{\isacharasterisk}{\kern0pt}{\isacharparenright}{\kern0pt}} defined in Isabelle. Prove the following
      property of it, i.e.\ that it is monotonically increasing in both arguments. You should prove
      it both, pen-and-paper and in Isabelle.%
\end{isamarkuptext}\isamarkuptrue%
\isacommand{lemma}\isamarkupfalse%
\ mult{\isacharunderscore}{\kern0pt}mono{\isacharcolon}{\kern0pt}\isanewline
\ \ \isakeyword{fixes}\ a\ b\ c\ d{\isacharcolon}{\kern0pt}{\isacharcolon}{\kern0pt}\ nat\ %
\isamarkupcmt{Note the alternative way of fixing the types of the constants%
}\isanewline
\ \ \isakeyword{assumes}\ {\isachardoublequoteopen}a\ {\isasymle}\ b{\isachardoublequoteclose}\ {\isachardoublequoteopen}c\ {\isasymle}\ d{\isachardoublequoteclose}\isanewline
\ \ \isakeyword{shows}\ {\isachardoublequoteopen}a\ {\isacharasterisk}{\kern0pt}\ c\ {\isasymle}\ b\ {\isacharasterisk}{\kern0pt}\ d{\isachardoublequoteclose}%
\begin{isamarkuptext}%
Your pen-and-paper proof should indicate

 1. what lemmas are you assuming,

 2. on what variable are you preforming induction,

 3. what are assumptions and the proof goals in the base case and the step case, and what is the
    induction hypothesis, and

 4. how are you instantiating the induction hypothesis, i.e. how do you show that its assumptions
    hold and which quantified variables are instantiated with which constants.

\paragraph{Hints:}
 1. The proof should be by induction

  2. If the induction hypothesis is too week, you might need to strengthen by generalising/
     universally quantifying over one of the variables. E.g. if you want to universally quantify on
     a variable \isa{x}, then you can use \isa{induction\ {\isachardot}{\kern0pt}{\isachardot}{\kern0pt}\ arbitrary{\isacharcolon}{\kern0pt}\ x}.

  3. The following lemmas suffice to prove the goal with only substitution, forward reasoning, and
     induction:

     \isa{refl{\isacharcomma}{\kern0pt}\ diff{\isacharunderscore}{\kern0pt}le{\isacharunderscore}{\kern0pt}self{\isacharcomma}{\kern0pt}\ mult{\isacharunderscore}{\kern0pt}le{\isacharunderscore}{\kern0pt}mono{\isadigit{2}}{\isacharcomma}{\kern0pt}\ mult{\isacharunderscore}{\kern0pt}zero{\isacharunderscore}{\kern0pt}right{\isacharcomma}{\kern0pt}\ le{\isadigit{0}}{\isacharcomma}{\kern0pt}\ mult{\isacharunderscore}{\kern0pt}Suc{\isacharunderscore}{\kern0pt}right{\isacharcomma}{\kern0pt}\ add{\isacharunderscore}{\kern0pt}mono{\isacharcomma}{\kern0pt}\ Suc{\isacharunderscore}{\kern0pt}eq{\isacharunderscore}{\kern0pt}plus{\isadigit{1}}{\isacharcomma}{\kern0pt}\ le{\isacharunderscore}{\kern0pt}diff{\isacharunderscore}{\kern0pt}conv}

  4. When you have a term like \isa{x\ {\isacharminus}{\kern0pt}\ y}, where \isa{x} is a \isa{nat}, in the goal or the assumptions, you 
     should perform proof by case analysis on whether \isa{y\ {\isasymle}\ x}. This is because of the way \isa{{\isacharminus}{\kern0pt}} is
     defined for natural numbers, where \isa{x\ {\isacharminus}{\kern0pt}\ y\ {\isacharequal}{\kern0pt}\ {\isadigit{0}}}, for any \isa{x\ {\isasymle}\ y}.%
\end{isamarkuptext}\isamarkuptrue%
%
\isadelimproof
\ \ %
\endisadelimproof
%
\isatagproof
\isacommand{using}\isamarkupfalse%
\ assms\isanewline
\isacommand{proof}\isamarkupfalse%
{\isacharparenleft}{\kern0pt}induction\ d\ arbitrary{\isacharcolon}{\kern0pt}\ c{\isacharparenright}{\kern0pt}\isanewline
\ \ \isacommand{case}\isamarkupfalse%
\ {\isadigit{0}}\isanewline
\ \ \isacommand{have}\isamarkupfalse%
\ {\isadigit{1}}{\isacharcolon}{\kern0pt}\ {\isachardoublequoteopen}c\ {\isacharequal}{\kern0pt}\ {\isadigit{0}}{\isachardoublequoteclose}\isanewline
\ \ \ \ \isacommand{using}\isamarkupfalse%
\ bot{\isacharunderscore}{\kern0pt}nat{\isacharunderscore}{\kern0pt}{\isadigit{0}}{\isachardot}{\kern0pt}extremum{\isacharunderscore}{\kern0pt}uniqueI{\isacharbrackleft}{\kern0pt}OF\ {\isachardoublequoteopen}{\isadigit{0}}{\isachardoublequoteclose}{\isacharparenleft}{\kern0pt}{\isadigit{2}}{\isacharparenright}{\kern0pt}{\isacharbrackright}{\kern0pt}\isanewline
\ \ \ \ \isacommand{{\isachardot}{\kern0pt}}\isamarkupfalse%
\isanewline
\ \ \isacommand{have}\isamarkupfalse%
\ {\isadigit{2}}{\isacharcolon}{\kern0pt}\ {\isachardoublequoteopen}a\ {\isacharasterisk}{\kern0pt}\ {\isadigit{0}}\ {\isacharequal}{\kern0pt}\ {\isadigit{0}}{\isachardoublequoteclose}\isanewline
\ \ \ \ \isacommand{using}\isamarkupfalse%
\ mult{\isacharunderscore}{\kern0pt}zero{\isacharunderscore}{\kern0pt}right\isanewline
\ \ \ \ \isacommand{{\isachardot}{\kern0pt}}\isamarkupfalse%
\ \isanewline
\isanewline
\ \ \isacommand{have}\isamarkupfalse%
\ {\isadigit{3}}{\isacharcolon}{\kern0pt}\ {\isachardoublequoteopen}a\ {\isacharasterisk}{\kern0pt}\ c\ {\isacharequal}{\kern0pt}\ {\isadigit{0}}{\isachardoublequoteclose}\isanewline
\ \ \ \ \isacommand{using}\isamarkupfalse%
\ {\isadigit{2}}\isanewline
\ \ \ \ \isacommand{apply}\isamarkupfalse%
{\isacharparenleft}{\kern0pt}subst\ {\isadigit{1}}{\isacharparenright}{\kern0pt}\isanewline
\ \ \ \ \isacommand{{\isachardot}{\kern0pt}}\isamarkupfalse%
\isanewline
\isanewline
\ \ \isacommand{have}\isamarkupfalse%
\ {\isadigit{4}}{\isacharcolon}{\kern0pt}\ {\isachardoublequoteopen}{\isadigit{0}}\ {\isasymle}\ b\ {\isacharasterisk}{\kern0pt}\ {\isadigit{0}}{\isachardoublequoteclose}\isanewline
\ \ \ \ \isacommand{using}\isamarkupfalse%
\ zero{\isacharunderscore}{\kern0pt}le\isanewline
\ \ \ \ \isacommand{{\isachardot}{\kern0pt}}\isamarkupfalse%
\isanewline
\isanewline
\ \ \isacommand{show}\isamarkupfalse%
\ {\isacharquery}{\kern0pt}case\isanewline
\ \ \ \ \isacommand{using}\isamarkupfalse%
\ {\isadigit{4}}\isanewline
\ \ \ \ \isacommand{apply}\isamarkupfalse%
{\isacharparenleft}{\kern0pt}subst\ {\isadigit{3}}{\isacharparenright}{\kern0pt}\isanewline
\ \ \ \ \isacommand{{\isachardot}{\kern0pt}}\isamarkupfalse%
\isanewline
\isanewline
\isacommand{next}\isamarkupfalse%
\isanewline
\ \ \isacommand{case}\isamarkupfalse%
\ {\isacharparenleft}{\kern0pt}Suc\ d{\isacharparenright}{\kern0pt}\isanewline
\ \ \isacommand{from}\isamarkupfalse%
\ {\isacartoucheopen}c\ {\isasymle}\ Suc\ d{\isacartoucheclose}\ \isacommand{have}\isamarkupfalse%
\ {\isadigit{1}}{\isacharcolon}{\kern0pt}\ {\isachardoublequoteopen}c\ {\isasymle}\ d\ {\isacharplus}{\kern0pt}\ {\isadigit{1}}{\isachardoublequoteclose}\isanewline
\ \ \ \ \isacommand{apply}\isamarkupfalse%
\ {\isacharparenleft}{\kern0pt}subst\ Suc{\isacharunderscore}{\kern0pt}eq{\isacharunderscore}{\kern0pt}plus{\isadigit{1}}{\isacharbrackleft}{\kern0pt}symmetric{\isacharbrackright}{\kern0pt}{\isacharparenright}{\kern0pt}\isanewline
\ \ \ \ \isacommand{{\isachardot}{\kern0pt}}\isamarkupfalse%
\isanewline
\ \ \isacommand{hence}\isamarkupfalse%
\ {\isachardoublequoteopen}c\ {\isacharminus}{\kern0pt}\ {\isadigit{1}}\ {\isasymle}\ d{\isachardoublequoteclose}\isanewline
\ \ \ \ \isacommand{apply}\isamarkupfalse%
{\isacharparenleft}{\kern0pt}subst\ le{\isacharunderscore}{\kern0pt}diff{\isacharunderscore}{\kern0pt}conv{\isacharparenright}{\kern0pt}\isanewline
\ \ \ \ \isacommand{{\isachardot}{\kern0pt}}\isamarkupfalse%
\isanewline
\isanewline
\ \ \isanewline
\isanewline
\ \ \isacommand{hence}\isamarkupfalse%
\ {\isadigit{2}}{\isacharcolon}{\kern0pt}\ {\isachardoublequoteopen}a\ {\isacharasterisk}{\kern0pt}\ {\isacharparenleft}{\kern0pt}c\ {\isacharminus}{\kern0pt}\ {\isadigit{1}}{\isacharparenright}{\kern0pt}\ {\isasymle}\ b\ {\isacharasterisk}{\kern0pt}\ d{\isachardoublequoteclose}\isanewline
\ \ \ \ \isanewline
\ \ \ \ \isacommand{apply}\isamarkupfalse%
{\isacharparenleft}{\kern0pt}rule\ Suc{\isachardot}{\kern0pt}IH{\isacharbrackleft}{\kern0pt}OF\ {\isacartoucheopen}a\ {\isasymle}\ b{\isacartoucheclose}{\isacharbrackright}{\kern0pt}{\isacharparenright}{\kern0pt}\isanewline
\ \ \ \ \isacommand{{\isachardot}{\kern0pt}}\isamarkupfalse%
\isanewline
\isanewline
\ \ \ \isanewline
\ \ \isanewline
\ \ \isacommand{show}\isamarkupfalse%
\ {\isacharquery}{\kern0pt}case\isanewline
\ \ \isacommand{proof}\isamarkupfalse%
{\isacharparenleft}{\kern0pt}cases\ c{\isacharparenright}{\kern0pt}\isanewline
\ \ \ \ \isacommand{case}\isamarkupfalse%
\ {\isadigit{0}}\isanewline
\ \ \ \ \isacommand{have}\isamarkupfalse%
\ {\isachardoublequoteopen}a\ {\isacharasterisk}{\kern0pt}\ c\ {\isacharequal}{\kern0pt}\ a\ {\isacharasterisk}{\kern0pt}\ {\isadigit{0}}{\isachardoublequoteclose}\isanewline
\ \ \ \ \ \ \isacommand{apply}\isamarkupfalse%
{\isacharparenleft}{\kern0pt}subst\ {\isadigit{0}}{\isacharparenright}{\kern0pt}\isanewline
\ \ \ \ \ \ \isacommand{{\isachardot}{\kern0pt}{\isachardot}{\kern0pt}}\isamarkupfalse%
\isanewline
\ \ \ \ \isacommand{also}\isamarkupfalse%
\ \isacommand{have}\isamarkupfalse%
\ {\isadigit{3}}{\isacharcolon}{\kern0pt}\ {\isachardoublequoteopen}{\isachardot}{\kern0pt}{\isachardot}{\kern0pt}{\isachardot}{\kern0pt}\ {\isacharequal}{\kern0pt}\ {\isadigit{0}}{\isachardoublequoteclose}\isanewline
\ \ \ \ \ \ \isacommand{apply}\isamarkupfalse%
{\isacharparenleft}{\kern0pt}subst\ mult{\isacharunderscore}{\kern0pt}zero{\isacharunderscore}{\kern0pt}right{\isacharparenright}{\kern0pt}\isanewline
\ \ \ \ \ \ \isacommand{{\isachardot}{\kern0pt}{\isachardot}{\kern0pt}}\isamarkupfalse%
\isanewline
\ \ \ \ \isacommand{also}\isamarkupfalse%
\ \isacommand{have}\isamarkupfalse%
\ {\isadigit{4}}{\isacharcolon}{\kern0pt}\ {\isachardoublequoteopen}{\isachardot}{\kern0pt}{\isachardot}{\kern0pt}{\isachardot}{\kern0pt}\ {\isasymle}\ b\ {\isacharasterisk}{\kern0pt}\ {\isacharparenleft}{\kern0pt}Suc\ d{\isacharparenright}{\kern0pt}{\isachardoublequoteclose}\isanewline
\ \ \ \ \ \ \isacommand{using}\isamarkupfalse%
\ le{\isadigit{0}}\isanewline
\ \ \ \ \ \ \isacommand{{\isachardot}{\kern0pt}}\isamarkupfalse%
\isanewline
\ \ \ \ \isacommand{finally}\isamarkupfalse%
\ \isacommand{show}\isamarkupfalse%
\ {\isacharquery}{\kern0pt}thesis\isanewline
\ \ \ \ \ \ \isacommand{{\isachardot}{\kern0pt}}\isamarkupfalse%
\isanewline
\ \ \isacommand{next}\isamarkupfalse%
\isanewline
\ \ \ \ \isacommand{case}\isamarkupfalse%
\ {\isacharparenleft}{\kern0pt}Suc\ nat{\isacharparenright}{\kern0pt}\isanewline
\ \ \ \ \isacommand{have}\isamarkupfalse%
\ {\isachardoublequoteopen}{\isacharparenleft}{\kern0pt}c\ {\isacharminus}{\kern0pt}\ {\isadigit{1}}{\isacharparenright}{\kern0pt}\ {\isacharequal}{\kern0pt}\ Suc\ nat\ {\isacharminus}{\kern0pt}{\isadigit{1}}{\isachardoublequoteclose}\isanewline
\ \ \ \ \ \ \isacommand{apply}\isamarkupfalse%
{\isacharparenleft}{\kern0pt}subst\ Suc{\isacharparenright}{\kern0pt}\isanewline
\ \ \ \ \ \ \isacommand{{\isachardot}{\kern0pt}{\isachardot}{\kern0pt}}\isamarkupfalse%
\isanewline
\ \ \ \ \isacommand{also}\isamarkupfalse%
\ \isacommand{have}\isamarkupfalse%
\ {\isachardoublequoteopen}{\isachardot}{\kern0pt}{\isachardot}{\kern0pt}{\isachardot}{\kern0pt}\ {\isacharequal}{\kern0pt}\ nat{\isachardoublequoteclose}\isanewline
\ \ \ \ \ \ \isacommand{apply}\isamarkupfalse%
{\isacharparenleft}{\kern0pt}subst\ diff{\isacharunderscore}{\kern0pt}Suc{\isacharunderscore}{\kern0pt}{\isadigit{1}}{\isacharparenright}{\kern0pt}\isanewline
\ \ \ \ \ \ \isacommand{{\isachardot}{\kern0pt}{\isachardot}{\kern0pt}}\isamarkupfalse%
\isanewline
\ \ \ \ \isacommand{finally}\isamarkupfalse%
\ \isacommand{have}\isamarkupfalse%
\ {\isachardoublequoteopen}c\ {\isacharminus}{\kern0pt}\ {\isadigit{1}}\ {\isacharequal}{\kern0pt}\ nat{\isachardoublequoteclose}\isanewline
\ \ \ \ \ \ \isacommand{{\isachardot}{\kern0pt}}\isamarkupfalse%
\isanewline
\isanewline
\ \ \ \ \isacommand{have}\isamarkupfalse%
\ {\isadigit{3}}{\isacharcolon}{\kern0pt}\ {\isachardoublequoteopen}a\ {\isacharasterisk}{\kern0pt}\ nat\ {\isasymle}\ b\ {\isacharasterisk}{\kern0pt}\ d{\isachardoublequoteclose}\isanewline
\ \ \ \ \ \ \isacommand{using}\isamarkupfalse%
\ {\isadigit{2}}\isanewline
\ \ \ \ \ \ \isacommand{apply}\isamarkupfalse%
{\isacharparenleft}{\kern0pt}subst\ {\isacartoucheopen}c\ {\isacharminus}{\kern0pt}\ {\isadigit{1}}\ {\isacharequal}{\kern0pt}\ nat{\isacartoucheclose}{\isacharbrackleft}{\kern0pt}symmetric{\isacharbrackright}{\kern0pt}{\isacharparenright}{\kern0pt}\isanewline
\ \ \ \ \ \ \isacommand{{\isachardot}{\kern0pt}}\isamarkupfalse%
\isanewline
\isanewline
\ \ \ \ \isacommand{have}\isamarkupfalse%
\ {\isachardoublequoteopen}a\ {\isacharasterisk}{\kern0pt}\ c\ {\isacharequal}{\kern0pt}\ a\ {\isacharasterisk}{\kern0pt}\ {\isacharparenleft}{\kern0pt}Suc\ nat{\isacharparenright}{\kern0pt}{\isachardoublequoteclose}\isanewline
\ \ \ \ \ \ \isacommand{apply}\isamarkupfalse%
{\isacharparenleft}{\kern0pt}subst\ \ Suc{\isacharparenright}{\kern0pt}\isanewline
\ \ \ \ \ \ \isacommand{{\isachardot}{\kern0pt}{\isachardot}{\kern0pt}}\isamarkupfalse%
\isanewline
\ \ \ \ \isacommand{also}\isamarkupfalse%
\ \isacommand{have}\isamarkupfalse%
\ {\isachardoublequoteopen}{\isachardot}{\kern0pt}{\isachardot}{\kern0pt}{\isachardot}{\kern0pt}\ \ {\isacharequal}{\kern0pt}\ a\ {\isacharplus}{\kern0pt}\ a\ {\isacharasterisk}{\kern0pt}\ nat{\isachardoublequoteclose}\isanewline
\ \ \ \ \ \ \isacommand{apply}\isamarkupfalse%
{\isacharparenleft}{\kern0pt}subst\ mult{\isacharunderscore}{\kern0pt}Suc{\isacharunderscore}{\kern0pt}right{\isacharparenright}{\kern0pt}\isanewline
\ \ \ \ \ \ \isacommand{{\isachardot}{\kern0pt}{\isachardot}{\kern0pt}}\isamarkupfalse%
\isanewline
\ \ \ \ \isacommand{finally}\isamarkupfalse%
\ \isacommand{have}\isamarkupfalse%
\ {\isadigit{4}}{\isacharcolon}{\kern0pt}\ {\isachardoublequoteopen}a\ {\isacharasterisk}{\kern0pt}\ c\ {\isacharequal}{\kern0pt}\ a\ {\isacharplus}{\kern0pt}\ a\ {\isacharasterisk}{\kern0pt}\ nat{\isachardoublequoteclose}\isanewline
\ \ \ \ \ \ \isacommand{{\isachardot}{\kern0pt}}\isamarkupfalse%
\isanewline
\isanewline
\ \ \ \ \isacommand{have}\isamarkupfalse%
\ {\isadigit{5}}{\isacharcolon}{\kern0pt}\ {\isachardoublequoteopen}b\ {\isacharasterisk}{\kern0pt}\ Suc\ d\ {\isacharequal}{\kern0pt}\ b\ {\isacharplus}{\kern0pt}\ b\ {\isacharasterisk}{\kern0pt}\ d{\isachardoublequoteclose}\isanewline
\ \ \ \ \ \ \isacommand{apply}\isamarkupfalse%
\ {\isacharparenleft}{\kern0pt}subst\ mult{\isacharunderscore}{\kern0pt}Suc{\isacharunderscore}{\kern0pt}right{\isacharparenright}{\kern0pt}\isanewline
\ \ \ \ \ \ \isacommand{{\isachardot}{\kern0pt}{\isachardot}{\kern0pt}}\isamarkupfalse%
\isanewline
\isanewline
\ \ \ \ \isacommand{have}\isamarkupfalse%
\ {\isachardoublequoteopen}a\ {\isacharplus}{\kern0pt}\ a\ {\isacharasterisk}{\kern0pt}\ nat\ {\isasymle}\ b\ {\isacharplus}{\kern0pt}\ b\ {\isacharasterisk}{\kern0pt}\ d{\isachardoublequoteclose}\isanewline
\ \ \ \ \ \ \isacommand{using}\isamarkupfalse%
\ add{\isacharunderscore}{\kern0pt}mono{\isacharbrackleft}{\kern0pt}OF\ {\isacartoucheopen}a\ {\isasymle}\ b{\isacartoucheclose}\ {\isadigit{3}}{\isacharbrackright}{\kern0pt}\isanewline
\ \ \ \ \ \ \isacommand{{\isachardot}{\kern0pt}}\isamarkupfalse%
\isanewline
\isanewline
\ \ \ \ \isacommand{then}\isamarkupfalse%
\ \isacommand{have}\isamarkupfalse%
\ {\isachardoublequoteopen}a\ {\isacharasterisk}{\kern0pt}\ c\ {\isasymle}\ b\ {\isacharplus}{\kern0pt}\ b\ {\isacharasterisk}{\kern0pt}\ d{\isachardoublequoteclose}\isanewline
\ \ \ \ \ \ \isacommand{apply}\isamarkupfalse%
{\isacharparenleft}{\kern0pt}subst\ {\isadigit{4}}{\isacharparenright}{\kern0pt}\isanewline
\ \ \ \ \ \ \isacommand{{\isachardot}{\kern0pt}}\isamarkupfalse%
\isanewline
\ \ \ \ \isacommand{then}\isamarkupfalse%
\ \isacommand{show}\isamarkupfalse%
\ {\isacharquery}{\kern0pt}thesis\isanewline
\ \ \ \ \ \ \isacommand{apply}\isamarkupfalse%
{\isacharparenleft}{\kern0pt}subst\ {\isadigit{5}}{\isacharparenright}{\kern0pt}\isanewline
\ \ \ \ \ \ \isacommand{{\isachardot}{\kern0pt}}\isamarkupfalse%
\isanewline
\ \ \isacommand{qed}\isamarkupfalse%
\isanewline
\isacommand{qed}\isamarkupfalse%
%
\endisatagproof
{\isafoldproof}%
%
\isadelimproof
%
\endisadelimproof
\ \isanewline
\isanewline
\isanewline
\isanewline
%
\isadelimtheory
\isanewline
%
\endisadelimtheory
%
\isatagtheory
\isacommand{end}\isamarkupfalse%
%
\endisatagtheory
{\isafoldtheory}%
%
\isadelimtheory
%
\endisadelimtheory
%
\end{isabellebody}%
\endinput
%:%file=Week_1_sol.tex%:%
%:%6=1%:%
%:%11=2%:%
%:%12=2%:%
%:%13=3%:%
%:%14=4%:%
%:%23=6%:%
%:%27=8%:%
%:%28=9%:%
%:%29=10%:%
%:%33=13%:%
%:%34=14%:%
%:%35=15%:%
%:%36=16%:%
%:%37=17%:%
%:%38=18%:%
%:%42=21%:%
%:%46=25%:%
%:%47=26%:%
%:%48=27%:%
%:%50=44%:%
%:%51=44%:%
%:%53=46%:%
%:%54=47%:%
%:%55=48%:%
%:%56=49%:%
%:%57=50%:%
%:%58=51%:%
%:%59=52%:%
%:%60=53%:%
%:%61=54%:%
%:%62=55%:%
%:%63=56%:%
%:%71=59%:%
%:%72=59%:%
%:%73=60%:%
%:%74=60%:%
%:%75=61%:%
%:%76=61%:%
%:%77=62%:%
%:%78=62%:%
%:%79=63%:%
%:%80=63%:%
%:%81=64%:%
%:%82=64%:%
%:%83=65%:%
%:%84=65%:%
%:%85=66%:%
%:%86=66%:%
%:%87=67%:%
%:%88=68%:%
%:%89=68%:%
%:%90=69%:%
%:%91=69%:%
%:%92=70%:%
%:%93=70%:%
%:%94=71%:%
%:%95=71%:%
%:%96=71%:%
%:%97=72%:%
%:%98=72%:%
%:%99=73%:%
%:%100=73%:%
%:%101=74%:%
%:%102=74%:%
%:%103=75%:%
%:%104=76%:%
%:%105=76%:%
%:%106=77%:%
%:%107=77%:%
%:%108=78%:%
%:%109=78%:%
%:%110=79%:%
%:%111=80%:%
%:%112=80%:%
%:%113=81%:%
%:%114=81%:%
%:%115=82%:%
%:%116=82%:%
%:%117=83%:%
%:%118=84%:%
%:%119=84%:%
%:%120=85%:%
%:%121=86%:%
%:%122=86%:%
%:%123=87%:%
%:%124=88%:%
%:%125=88%:%
%:%126=89%:%
%:%127=89%:%
%:%128=90%:%
%:%129=91%:%
%:%130=91%:%
%:%131=92%:%
%:%132=92%:%
%:%133=93%:%
%:%134=94%:%
%:%135=94%:%
%:%136=95%:%
%:%137=95%:%
%:%138=96%:%
%:%139=97%:%
%:%149=104%:%
%:%153=106%:%
%:%154=107%:%
%:%155=108%:%
%:%157=131%:%
%:%158=131%:%
%:%159=132%:%
%:%160=132%:%
%:%161=132%:%
%:%162=133%:%
%:%163=134%:%
%:%165=137%:%
%:%166=138%:%
%:%167=139%:%
%:%168=140%:%
%:%169=141%:%
%:%170=142%:%
%:%171=143%:%
%:%172=144%:%
%:%173=145%:%
%:%174=146%:%
%:%175=147%:%
%:%176=148%:%
%:%177=149%:%
%:%178=150%:%
%:%179=151%:%
%:%180=152%:%
%:%181=153%:%
%:%182=154%:%
%:%183=155%:%
%:%184=156%:%
%:%185=157%:%
%:%186=158%:%
%:%187=159%:%
%:%187=160%:%
%:%188=161%:%
%:%189=162%:%
%:%190=163%:%
%:%191=164%:%
%:%195=170%:%
%:%199=170%:%
%:%200=170%:%
%:%201=171%:%
%:%202=171%:%
%:%203=172%:%
%:%204=172%:%
%:%205=173%:%
%:%206=173%:%
%:%207=174%:%
%:%208=174%:%
%:%209=175%:%
%:%210=175%:%
%:%211=176%:%
%:%212=176%:%
%:%213=177%:%
%:%214=177%:%
%:%215=178%:%
%:%216=178%:%
%:%217=179%:%
%:%218=180%:%
%:%219=180%:%
%:%220=181%:%
%:%221=181%:%
%:%222=182%:%
%:%223=182%:%
%:%224=183%:%
%:%225=183%:%
%:%226=184%:%
%:%227=185%:%
%:%228=185%:%
%:%229=186%:%
%:%230=186%:%
%:%231=187%:%
%:%232=187%:%
%:%233=188%:%
%:%234=189%:%
%:%235=189%:%
%:%236=190%:%
%:%237=190%:%
%:%238=191%:%
%:%239=191%:%
%:%240=192%:%
%:%241=192%:%
%:%242=193%:%
%:%243=194%:%
%:%244=194%:%
%:%245=195%:%
%:%246=195%:%
%:%247=196%:%
%:%248=196%:%
%:%249=196%:%
%:%250=197%:%
%:%251=197%:%
%:%252=198%:%
%:%253=198%:%
%:%254=199%:%
%:%255=199%:%
%:%256=200%:%
%:%257=200%:%
%:%258=201%:%
%:%259=201%:%
%:%260=202%:%
%:%261=203%:%
%:%262=204%:%
%:%263=205%:%
%:%264=205%:%
%:%265=206%:%
%:%266=207%:%
%:%267=207%:%
%:%268=208%:%
%:%269=208%:%
%:%270=209%:%
%:%271=210%:%
%:%272=211%:%
%:%273=212%:%
%:%274=212%:%
%:%275=213%:%
%:%276=213%:%
%:%277=214%:%
%:%278=214%:%
%:%279=215%:%
%:%280=215%:%
%:%281=216%:%
%:%282=216%:%
%:%283=217%:%
%:%284=217%:%
%:%285=218%:%
%:%286=218%:%
%:%287=218%:%
%:%288=219%:%
%:%289=219%:%
%:%290=220%:%
%:%291=220%:%
%:%292=221%:%
%:%293=221%:%
%:%294=221%:%
%:%295=222%:%
%:%296=222%:%
%:%297=223%:%
%:%298=223%:%
%:%299=224%:%
%:%300=224%:%
%:%301=224%:%
%:%302=225%:%
%:%303=225%:%
%:%304=226%:%
%:%305=226%:%
%:%306=227%:%
%:%307=227%:%
%:%308=228%:%
%:%309=228%:%
%:%310=229%:%
%:%311=229%:%
%:%312=230%:%
%:%313=230%:%
%:%314=231%:%
%:%315=231%:%
%:%316=231%:%
%:%317=232%:%
%:%318=232%:%
%:%319=233%:%
%:%320=233%:%
%:%321=234%:%
%:%322=234%:%
%:%323=234%:%
%:%324=235%:%
%:%325=235%:%
%:%326=236%:%
%:%327=237%:%
%:%328=237%:%
%:%329=238%:%
%:%330=238%:%
%:%331=239%:%
%:%332=239%:%
%:%333=240%:%
%:%334=240%:%
%:%335=241%:%
%:%336=242%:%
%:%337=242%:%
%:%338=243%:%
%:%339=243%:%
%:%340=244%:%
%:%341=244%:%
%:%342=245%:%
%:%343=245%:%
%:%344=245%:%
%:%345=246%:%
%:%346=246%:%
%:%347=247%:%
%:%348=247%:%
%:%349=248%:%
%:%350=248%:%
%:%351=248%:%
%:%352=249%:%
%:%353=249%:%
%:%354=250%:%
%:%355=251%:%
%:%356=251%:%
%:%357=252%:%
%:%358=252%:%
%:%359=253%:%
%:%360=253%:%
%:%361=254%:%
%:%362=255%:%
%:%363=255%:%
%:%364=256%:%
%:%365=256%:%
%:%366=257%:%
%:%367=257%:%
%:%368=258%:%
%:%369=259%:%
%:%370=259%:%
%:%371=259%:%
%:%372=260%:%
%:%373=260%:%
%:%374=261%:%
%:%375=261%:%
%:%376=262%:%
%:%377=262%:%
%:%378=262%:%
%:%379=263%:%
%:%380=263%:%
%:%381=264%:%
%:%382=264%:%
%:%383=265%:%
%:%384=265%:%
%:%385=266%:%
%:%393=266%:%
%:%394=267%:%
%:%395=268%:%
%:%396=269%:%
%:%399=270%:%
%:%404=271%:%
